\subsection{The cubic descent for all multipliers}

\begin{statement}{D23.F23}{Proved}{The whole gentle-group descent and its relative quadratic character}
\end{statement}
The endomorphism $\mathcal D_E$ of \textup{D23.F22} extends canonically
to a continuous endomorphism of the full actual gentle group
$\mathcal G=\widehat{GT}_0$, preserving its multiplier.
For $g=(\lambda,f)$ put
\[
 m=(\lambda-1)/2,\quad
 \varepsilon\in\{1,-1\},\quad\varepsilon\equiv\lambda\pmod3,
 \quad b=(1-\varepsilon)/2,
 \qquad C_g=r^{-m}f(x,r).
\]
Define
\[
 \psi_g(r)=r^\lambda,\qquad
 \psi_g(s)=f(r,srs)^{-1}s,\qquad
 \theta_g=\operatorname{Inn}(\tau^b C_g^{-1})\psi_g.
\]
Then $f_{0,g}=s\theta_g(s)$ belongs to $F$ and has abelianization
$(u_g,-u_g)$. The full descent is the explicit normalized automorphism
\[
 \boxed{\quad
 \mathcal D_E(g)=\operatorname{Inn}(x^{u_g})\theta_g,
 \quad f_{\mathcal D_E(g)}=y^{u_g}f_{0,g}x^{-u_g}\in F'.
 \quad}                                                     \tag{D23.F23.1}
\]
This construction is independent of any Galois lift of the multiplier.

Using \textup{D23.P13}, put $\gamma=\Delta\circ\mathcal D_E$.
Its restriction to $\mathcal H=\ker\lambda$ is the duplication
defect $d_2=a_2+\Psi_3/2$. There is a unique continuous character
$\chi:\widehat{\mathbb Z}^{\times}\to\mathbb F_2$ such that
$\gamma|_{G_{\mathbb Q}}=\chi\circ\lambda$.
Consequently
\[
 \widehat d_2=\gamma+\chi\circ\lambda,\qquad
 \widehat e_3=\Delta+\widehat d_2                             \tag{D23.F23.2}
\]
are canonical continuous relative quadratic characters of
$\mathcal G$, restricting respectively to
\[
 a_2+\Psi_3/2,\qquad a_3+\Psi_3/2
\]
on $\mathcal H$. In particular
$\Delta=\widehat d_2+\widehat e_3$ on the full gentle group.
This does not assert nontriviality or independence of these two
characters, and does not determine the whole relative character space.
It also does not assert that the uncorrected $\gamma$ vanishes on
all of $G_{\mathbb Q}$.
By uniqueness of relative extension, $\widehat d_2$ is the
character $d$ already constructed in D23.C1. The new assertion
is its realization through the whole-group cubic descent.

\begin{proof}
As in \textup{D23.F22}, quotient the input free profinite group by
the closed normal subgroup generated by $X^3,Y^3$. The quotient
is $E=C_3\amalg C_3$, with $X=r,Y=srs$. Both the actual input
automorphism and its inverse preserve this normal subgroup, since
$\lambda$ is a unit. They therefore induce an automorphism of $E$.
Write $G=f(r,srs)$. The two-cycle gives $sGs=G^{-1}$, so
$(G^{-1}s)^2=1$. The formulas for $\psi_g$ define an endomorphism
of $B$, whose restriction to $E$ is the induced automorphism and
whose action on $B/E=C_2$ is the identity. It is an automorphism.
The normalized pair-composition formula proves directly that
$g\mapsto\psi_g$ is a homomorphism, including arbitrary multipliers.

For clarity we derive the general peripheral conjugator. Write
$Z=(XY)^{-1}$, $G_0=f(Z,X)$, $H_0=f(Y,Z)$ and
$m=(\lambda-1)/2$. The three-cycle for $(\lambda,f)$ is
\[
 G_0 Z^m H_0 Y^m f X^m=1.
\]
Right composition with the actual element $c=(-1,1)$ gives the
pair $(-\lambda,f)$ and its three-cycle
\[
 G_0 Z^{-m-1}H_0Y^{-m-1}fX^{-m-1}=1.
\]
Eliminating $H_0$ from these equations gives
\[
 G_0Z^{-\lambda}G_0^{-1}
   =X^{m+1}f^{-1}Y^\lambda fX^m
   =X^{-m}\Phi_g(Z^{-1})X^m.
\]
Thus
\[
 \Phi_g(Z)=X^m f(Z,X)Z^\lambda f(Z,X)^{-1}X^{-m}.             \tag{D23.F23.3}
\]
Evaluate this at $(r,srs)$, where $Z=y=(rs)^{-2}$.
Since $x=r^{-1}yr$ and $\psi_g(r)=r^\lambda$, we obtain
\[
 \psi_g(x)=r^{m-\lambda}f(y,r)y^\lambda f(y,r)^{-1}r^{\lambda-m}
          =C_g x^\lambda C_g^{-1}.
\]
In the last equality use $m-\lambda+1=-m$ and
$f(y,r)r=rf(x,r)$.

The element $f(x,r)$ is pure, and $\psi_g$ induces on $S_3$
the automorphism $s\mapsto s$, $r\mapsto r^\varepsilon$.
Consequently $v=C_g^{-1}\psi_g(\tau)C_g$ has the same
permutation as $\tau^\lambda=\tau\bmod F$ and satisfies
$v^2=\tau^{2\lambda}$.
Both $v$ and $\tau^\lambda$ commute with $x^\lambda$.
Since $\lambda$ is a unit, they commute with $x$ as well.
Therefore $v\tau^{-\lambda}$ lies in
$C_F(x)=\overline{\langle x\rangle}$.
The equality of squares and injectivity of multiplication by two
in $\widehat{\mathbb Z}$ show that it is one. We have proved
\[
 \psi_g(\tau)=C_g\tau^\lambda C_g^{-1}.
                                                               \tag{D23.F23.4}
\]

The correction by $\tau^b$ is essential for the permutation
normalization. The two possible cases are
\[
\begin{array}{c|c|c|c|c}
 \lambda\bmod6&\varepsilon&m\bmod3&b&
       \tau^b C_g^{-1}\bmod F\\\hline
 1&1&0&0&1\\
 5&-1&2&1&(sr)r^2=s.
\end{array}
\]
In the first case $\psi_g$ already acts identically on $S_3$.
In the second it fixes $s$ and inverts $r$, and conjugation by
$s$ changes this into the identity on $S_3$. Thus in both cases
$\theta_g$ induces the identity on $S_3$; by (D23.F23.4) it
sends $\tau$ to $\tau^\lambda$.

It follows that $f_{0,g}=s\theta_g(s)$ belongs to $F$.
The order-two relation $\theta_g(s)^2=1$ gives the exact identity
$s f_{0,g}s=f_{0,g}^{-1}$. On $F^{\rm ab}$, conjugation by $s$
exchanges the two coordinates. Hence if
$f_{0,g}^{\rm ab}=(u,v)$, then $(v,u)=(-u,-v)$, or $v=-u$.
This proves the asserted abelian form. Conjugation by $x^{u_g}$
changes the pair coordinate to $y^{u_g}f_{0,g}x^{-u_g}$, whose
abelianization is zero, and preserves the cusp formula.

The normalized automorphism restricts to the actual free-group
automorphism
$x\mapsto x^\lambda$, $y\mapsto f_E^{-1}y^\lambda f_E$.
Its two-cycle follows from the order-two relation just used.
Its three-cycle follows from the order-three relation in $B$:
its image of $r$ is $f_E^{-1}r x^{-m}$, and expansion of its
cube, using the cyclic action of $r$ on $x,y,(xy)^{-1}$, gives
\[
 f_E(z,x)z^m f_E(y,z)y^m f_E(x,y)x^m=1.
\]
Thus (D23.F23.1) is an actual element of $\mathcal G$ of
multiplier $\lambda$.

For the composition law, each (D23.F23.1) is an inner correction
of $\psi_g$. Since $\psi$ is a homomorphism, its normalized
corrections have the same outer composition law. Two such
corrections with the same multiplier both induce the identity
on $S_3$, so an inner element relating them has trivial image
in $S_3$, whose center is trivial. It is therefore in $F$.
Since both automorphisms send $\tau$ to $\tau^\lambda$, that
element centralizes $x$ and is $x^n$. Conjugation by $x^n$
changes the pair coordinate by the abelian vector $(-n,n)$.
Both coordinates belong to $F'$, so $n=0$. This proves uniqueness
of the normalization and hence the homomorphism law.
All operations, the locally constant choice of $b$, and extraction
of $u_g$ from $F^{\rm ab}$ are continuous, proving continuity.
At $\lambda=1$ we have $b=m=0$ and, by \textup{D23.F22},
$u_g=-3\Psi_3(g)/2$. Thus the extension restricts to that of
\textup{D23.F22}.

Now $\gamma=\Delta\mathcal D_E$ is a continuous character on
$\mathcal G$. Statements \textup{D23.F22} and \textup{D23.P13}
give $\gamma|_{\mathcal H}=d_2$.
The arithmetic comparison needed here is precisely
\textup{DELTA\_ON\_GT A.103}, final paragraph after
\textup{(MD.4e)}: on $G_{\mathbb Q}\cap\mathcal H$,
Anderson duplication in the specified intrinsic convention gives
$d_2=0$. Equivalently this is the physical-descent arithmetic
compatibility of \textup{A.62}, combined with the universal
half-character formula \textup{A.107 (MD.8a)}.
We use this supplied arithmetic comparison, and do not replace
it by an unmarked assertion about a Kummer exponent.

Therefore $\gamma|_{G_{\mathbb Q}}$ vanishes on the kernel of
the cyclotomic character. The cyclotomic map
$G_{\mathbb Q}\twoheadrightarrow\widehat{\mathbb Z}^{\times}$
is a continuous surjection of compact profinite groups, hence
its target is the topological quotient by that kernel.
There is consequently a unique continuous character $\chi$
with the stated factorization. The formulas (D23.F23.2) define
continuous characters and show directly that they vanish on
$G_{\mathbb Q}$. Their restrictions to $\mathcal H$ are as
claimed. Their uniqueness as relative extensions of those
restrictions follows from
$\mathcal G=\mathcal H G_{\mathbb Q}$, the same decomposition
used in \textup{LM7\_PUSH E.65}. This also proves the last
displayed identity on all of $\mathcal G$.
\end{proof}

\paragraph{Geometric normalization.}
For the cubic eigenlocus put
$n=(1+\sqrt{-3})/2$, $\bar n=(1-\sqrt{-3})/2$ and
\[
 t=\frac{\lambda-\bar n}{\lambda-n},\qquad b=1-t^3.
\]
Direct substitution in the Legendre expression gives
\[
 \frac{256(1-\lambda+\lambda^2)^3}
      {1728\lambda^2(1-\lambda)^2}
       =\frac{4(b-1)}{b^2}.                                \tag{D23.F23.5}
\]
The pole values $n,\bar n$ correspond to $b=\infty,1$;
the three cusps $\lambda=0,1,\infty$ correspond to $b=0$.
The degree-two map on the right is the harmonic map
$4u(1-u)$ after $u=1/b$. This records the precise relation
between the $C_3\amalg C_3$ and $C_2\amalg C_3$ orbifolds.
It does not by itself identify the two GT actions: the cyclic
distribution filling retains $0,1,\infty$ in the $t$-sphere,
whereas the cubic modular filling removes the two poles and
retains the three roots of unity. The functorial comparison for
these different fillings remains the required additional step.
